\section{Discussion}
\label{sec:discussion}

\paragraph{Transformation subtypes in the context of established LUAD biology.}
The transcriptional signatures of the three subtypes are broadly consistent with the
state-based LUAD classification from TCGA \cite{tcga2014luad}.
Subtype~0, which is enriched for \textit{EGFR} mutations and shows suppressed
cell-cycle activation, resembles the Terminal Respiratory Unit (TRU) subtype,
characterised by a differentiated transcriptional state and \textit{EGFR}-driven
oncogenesis.
Subtype~1, with its dominant E2F and G2-M programmes, parallels the Proximal
Proliferative (PP) subtype.
Subtype~2, defined by interferon signalling activation, is broadly consistent with the
Proximal Inflammatory (PI) subtype, although the co-occurrence of E2F upregulation in
subtype~2 does not map cleanly onto classical PI biology.
Importantly, state-based and trajectory-based subtypes are conceptually distinct: TRU,
PP, and PI describe what a tumour looks like, whereas the MOSAIC subtypes describe
how the patient got there.
The parallels are therefore informative rather than expected, and their correspondence
warrants formal testing in a larger dataset.

The mutation enrichment result provides an independent line of evidence for the
biological coherence of the subtype structure.
Because \textit{EGFR} mutations and ALK rearrangements are mutually exclusive oncogenic
events in LUAD, the observation that both are differentially distributed within the same
subtype (enriched and depleted, respectively, in subtype~0) implies that the clustering
captures biologically meaningful patient groupings rather than statistical artefacts.

\paragraph{The L2 normalisation choice and cluster geometry.}
Clustering was performed on L2-normalised $\Delta Z$, which focuses on transformation
\emph{direction} rather than magnitude.
The practical consequence is that two patients with very different transformation
extents but qualitatively similar factor profiles are assigned to the same subtype,
while magnitude is treated as an independent dimension.
This choice is supported by the absence of a significant association between
transformation magnitude and tumour stage, which argues against magnitude being a
simple proxy for disease state.
The resulting silhouette scores ($\approx 0.18$) are modest, reflecting a partially
continuous manifold rather than discrete clusters; this is expected for a cohort of
100 patients stratified by a 13-dimensional biological signal.

\paragraph{Limitations.}
The most consequential limitation is sample size.
With $n = 102$ patients, many of whom were censored before the study endpoint, the
study is underpowered for survival analysis and the null log-rank result is
uninformative rather than negative evidence.
Biological interpretation relies exclusively on transcriptomic data at two stages:
GSEA on MOFA factor loadings uses RNA weights only, and cluster differential expression
uses tumour-minus-normal RNA rather than the multi-omics signal used in the MOFA fit.
The proteomic and phosphoproteomic modalities therefore contribute to the latent space
but are currently absent from the interpretation layer, leaving a portion of the
multi-omics information unexploited.
Finally, all results are derived from a single cohort; in the absence of external
validation, the reported subtypes should be treated as hypotheses rather than
established molecular entities.

\paragraph{Future directions.}
Multi-cancer validation using other CPTAC cohorts (e.g.\ breast or colorectal) would
test whether latent trajectory structure is a general property of matched multi-omics
tumourigenesis or specific to LUAD.
Kinase activity inference from phosphoproteomics factor loadings (e.g.\ using
kinase-substrate databases such as PhosphoSitePlus) would bring the third omics
modality into the interpretation layer and link transformation subtypes to specific
regulatory programmes.
Fitting a continuous manifold model on $\Delta Z$ vectors, such as diffusion
pseudotime, would complement the discrete clustering by revealing whether the three
subtypes represent distinct attractors or points along a shared trajectory.
The stochastic nature of MOFA fitting also deserves attention: comparing factor
structures across models fitted with different random seeds would clarify which
factors are stable properties of the data and which are more sensitive to
initialisation.
Finally, the current analysis uses MAD-based preselection on three omics views.
Expanding the feature selection strategy or incorporating additional views (e.g.\
copy-number alterations or miRNA) may better separate the transformation modes and
reduce the current dependence on a single, view-agnostic ranking criterion.
